\documentclass[11pt]{article}
\usepackage[margin=1in]{geometry}
\usepackage{amsmath,amssymb,amsthm}
\usepackage[T1]{fontenc}
\usepackage{hyperref}
\newcommand{\NE}{\operatorname{NE}}
\newcommand{\Stab}{\operatorname{Stab}_0}
\newcommand{\Reg}{\operatorname{Reg}}
\newcommand{\SwapReg}{\operatorname{SwapReg}}
\newcommand{\pos}[1]{[#1]_+}
\newtheorem{proposition}{Proposition}
\title{Strategic axioms selecting stable sets themselves}
\author{}
\date{Status checked: 17 October 2026}
\begin{document}
\maketitle
\section*{Status}
This is an explicitly labelled \emph{research programme}, not a fixed
theorem-or-counterexample conjecture. No proposed new axiom or solution
is supplied. It accompanies, rather than duplicates, the containment
question S1.
\section*{Objects and conventions}
A finite normal-form game $G=(I,(A_i),(u_i))$ has mixed-strategy space
$\Sigma=\prod_{i\in I}\Delta(A_i)$ and Nash set $\NE(G)$.
A solution concept is a rule $\Phi$ defined on \emph{all} finite
normal-form games, including games with additional players and actions:
\[
 \varnothing\ne\Phi(G)\subseteq
 \{K\subseteq\NE(G):K\ne\varnothing,\ K\text{ closed and connected}\}.
\]
Write $\Stab(G)$ for Mertens $0$-stable sets, with rational cohomology,
using exactly Definition 2.4 of the version of Govindan--Wilson cited
below. The zero specifies the coefficient characteristic, not an
equilibrium's local index. To fix the perturbation model, for
$0<\varepsilon\leq1$ set
\[
 P_\varepsilon=\{\delta\tau:0\leq\delta\leq\varepsilon,\ \tau\in\Sigma\},
 \quad \Sigma_i(\eta)=\{x_i\in\Delta(A_i):x_i(a)\geq\eta_i(a)\ \forall a\in A_i\},
\]
where $\eta\in P_1$. Let $G(\eta)$ be $G$ with each player's
strategy space constrained to $\Sigma_i(\eta)$, and let
$E_G=\{(\eta,x)\in P_1\times\Sigma:x\in\NE(G(\eta))\}$,
equipped with the projection $(\eta,x)\mapsto\eta$.
The definition combines cohomological essentiality
of this projection with the source's connexity condition. Definition 2.4,
including its boundary conventions, is normative here.

Use the exact cross-game axioms $B,D^*,I^*$ of that same version,
Sections 3--5 (Definitions 4.4 and 5.2), rather than informal substitutes.
Axiom $B$ requires every selected set to contain a profile equivalent to
a quasi-perfect equilibrium in every perfect-recall extensive-form
representation. Axiom $D^*$ concerns additions of strongly irrelevant
dominated strategies. Axiom $I^*$ concerns strong embeddings, including
additional players; its comparisons use the source's projection maps.
The full technical restrictions in those definitions are part of the
problem. Arbitrary dominated-action additions or arbitrary embeddings
are not being substituted for them.

\section*{Programme S2}
Find and state a precise additional axiom $\mathcal A$ such that
\[
 \Stab\text{ satisfies }\mathcal A,
 \qquad
 B+D^*+I^*+\mathcal A
    \ \Longrightarrow\ 
 \forall G\quad\Phi(G)\subseteq\Stab(G).
\]
The preferred domain is all finite games. A generic-game result is a
clearly labelled intermediate target. Each such intermediate proposal
must specify its representation, payoff parameter space, exceptional
set, and game domain. The axiom should have an
independent interpretation in strategic choice, beliefs, or invariance.
An axiom that simply says the selected sets are stable would restate the
desired conclusion and would not address that research motivation.

\section*{Deliverables and scope}
A proposed answer must specify $\mathcal A$ mathematically, its game
domain, the interpretation motivating it, a proof that $0$-stability
satisfies it, and a proof or counterexample for the implication.
Once a particular $\mathcal A$ is fixed, that implication is a definite
mathematical question. The phrase ``independently motivated'' is a
criterion for evaluating proposals, not a hidden formal predicate.

This programme asks for $\Phi(G)\subseteq\Stab(G)$, not equality.
Equality would additionally require every stable set to be selected.
The condition that stability itself satisfies the axioms prevents an
empty or inconsistent axiom system from counting as a solution.

\section*{Existing results and editorial handling}
The generic containment theorem is already known. Govindan--Wilson
discuss minimality as a way to strengthen selection, and explain why
minimality alone does not supply their desired decision-theoretic
rationale \cite[Sections 6.1 and 10]{gw}. Thus minimality must not be
presented here as a newly discovered solution of this programme.
No exact we-solve target was found. If a problem entry must be a fixed
mathematical conjecture, retain this document as motivation for S1
instead of opening an independent problem entry. The submission draft
explicitly preserves that limitation.

\section*{Preparation}
The original research map was generated using OpenAI GPT 6 Astra Pro.
Source checks, formulation, and LaTeX preparation used a Codex assistant
based on GPT-6. Exact model versions and reasoning settings are unknown.
The cited researchers retain credit for their results. Source statements
and scope were checked; no complete proof audit or novelty certification
is claimed.
\paragraph{Origin of the questions.}
The open-problem suggestions motivating this entry are credited to
Rida Laraki's Facebook comment~\cite{laraki-facebook}.
\begin{thebibliography}{9}
\bibitem{gw}
Srihari Govindan and Robert Wilson.
\emph{Axiomatic Equilibrium Selection in Generic Extensive-Form Games}.
arXiv:2504.16908v2, 26 August 2026.
Definitions 2.4, 2.5, 4.4, 5.2; Axioms 3.1, 4.7, 5.4; Theorem 6.1; Sections 6.1 and 10.
\url{https://arxiv.org/html/2504.16908v2}.
\bibitem{laraki-facebook}
Rida Laraki.
\emph{Facebook comment suggesting open problems in game theory}.
\end{thebibliography}
\end{document}
